\documentclass[a4paper,12pt]{article}
\usepackage[utf8]{inputenc}
\usepackage[main=german,provide=*]{babel}
\usepackage{amsmath}
\usepackage{amssymb}
\usepackage{amsthm}
\usepackage{geometry}
\usepackage{graphicx}
\usepackage[hidelinks]{hyperref}
\usepackage{listings}
\usepackage{xcolor}
\usepackage{enumitem}
\usepackage{microtype}

\geometry{margin=2.5cm}
\setlength{\emergencystretch}{3em}
\sloppy

\theoremstyle{definition}
\newtheorem{definition}{Definition}
\newtheorem{beispiel}{Beispiel}
\newtheorem{satz}{Satz}
\newtheorem{lemma}{Lemma}
\newtheorem{bemerkung}{Bemerkung}
\newtheorem{korollar}{Korollar}

\setlist[itemize,1]{label=--}

% ---------- Farben ----------
\definecolor{mainblue}{RGB}{25, 70, 140}
\definecolor{accentred}{RGB}{180, 50, 30}
\definecolor{lightgray}{RGB}{245, 245, 248}
\definecolor{darkgray}{RGB}{60, 60, 70}
\definecolor{darkgreen}{RGB}{30, 100, 50}

\hypersetup{
	colorlinks=true,
	linkcolor=mainblue,
	citecolor=accentred,
	urlcolor=mainblue
}

%  Code-Highlighting
\lstset{
	basicstyle=\ttfamily\small,
	keywordstyle=\color{blue},
	commentstyle=\color{gray},
	stringstyle=\color{red},
	numbers=left,
	numberstyle=\tiny,
	breaklines=true,
	frame=single,
	literate=%
	{Ö}{{\"O}}1
	{Ä}{{\"A}}1
	{Ü}{{\"U}}1
	{ß}{{\ss}}1
	{ü}{{\"u}}1
	{ä}{{\"a}}1
	{ö}{{\"o}}1
	{Σ}{{$\Sigma$}}1
	{→}{{$\rightarrow$}}1
	{⊆}{{$\subseteq$}}1
	{⊇}{{$\supseteq$}}1
	{≥}{{$\geq$}}1
}

\title{Der paradoxe Gewinn im Reziproken:\\
Die Rolle des Exponenten $n$ zur Basis $2$\\
in der Signaturfunktion des Subgraph Algorithmus}
\author{Stephan Epp}
\date{11. Juni 2026}

\begin{document}

\maketitle

\tableofcontents
\newpage

\section{Einführung}

Der Subgraph Algorithmus vergleicht zwei Graphen $G$ und $G'$ mittels eindeutiger
Spalten-Signaturen und zyklischer Rotation und erreicht dabei eine optimale
Laufzeit von $O(n^3)$. Zentrales Element der Konstruktion ist die Signaturfunktion
\[
	\sigma_j = \sum_{i=0}^{n-1} A_{ij} \cdot 2^i + j \cdot 2^n,
\]
in welcher der Exponent $n$ zur Basis $2$ gleich zweifach auftritt: einmal als
obere Grenze der Zeilensumme (\emph{Zeilenkomponente}, $\sum_i A_{ij}2^i < 2^n$)
und einmal als multiplikatives Gewicht der Spalte (\emph{Spaltenkomponente},
$j \cdot 2^n$).

Die vorliegende Arbeit untersucht ein auf den ersten Blick paradoxes Phänomen,
das wir als \emph{paradoxen Gewinn im Reziproken} bezeichnen: Während der
Term $2^n$ mit wachsendem $n$ exponentiell divergiert, konvergiert sein
Reziprokes $2^{-n}$ exponentiell gegen $0$. Genau diese Divergenz-Konvergenz-Dualität
ist jedoch keine Schwäche, sondern die strukturelle Ursache dafür, dass
\begin{itemize}
	\item die Signaturfunktion $\sigma$ \emph{injektiv} ist (Eindeutigkeit, vgl.
	      Lemma~1 der Originalarbeit \cite{epp2026subgraph}),
	\item Zeilen- und Spalteninformation durch eine einzige reziproke Skalierung
	      $\sigma_j \cdot 2^{-n}$ \emph{verlustfrei und ohne zusätzlichen Speicheraufwand}
	      getrennt werden können,
	\item und diese Trennung exakt die Struktur ist, welche die untere
	      Schranke $\Omega(n^3)$ für das zyklische Subgraph-Problem erzwingt.
\end{itemize}

\subsection{Problemstellung}

Sei $n \in \mathbb{N}$, $n \geq 1$, und sei $A \in \{0,1\}^{n \times n}$ eine
Adjazenzmatrix. Für jede Spalte $j \in \{0,\ldots,n-1\}$ sei
\[
	r_j := \sum_{i=0}^{n-1} A_{ij} \cdot 2^i \in \{0, 1, \ldots, 2^n - 1\}
\]
die \emph{Zeilenkomponente} und $\sigma_j = r_j + j \cdot 2^n$ die zugehörige
Signatur.

\textbf{Frage:} Welche formalen Eigenschaften besitzt die Abbildung
$\sigma_j \mapsto (j, r_j)$, und in welchem Sinn stellt der Ãœbergang
$2^n \mapsto 2^{-n}$ einen \emph{Gewinn} dar, obwohl er auf den ersten Blick
nur eine triviale algebraische Inversion ist?

\subsection{Motivation}

In der Praxis erscheinen exponentiell wachsende Größen wie $2^n$ zunächst
als Nachteil: Speicherplatz und Zahlendarstellung wachsen mit $n$. Das
\emph{Reziproke} $2^{-n}$ hingegen erscheint als vernachlässigbar kleine Größe,
die „verschwindet''. Der paradoxe Gewinn besteht darin, dass die
Multiplikation mit $2^{-n}$ \emph{nicht} Information vernichtet, sondern
sie in zwei disjunkte, unmittelbar auslesbare Komponenten – Ganzzahlanteil
(Spaltenindex $j$) und Fraktionalanteil (Zeileninformation $r_j / 2^n$) –
zerlegt. Dieser Mechanismus ist die mathematische Grundlage dafür, dass der
Subgraph Algorithmus auf zusätzliche Indexstrukturen verzichten kann und
dennoch $O(n^3)$ erreicht.

\section{Grundlagen}

\subsection{Definitionen}

\begin{definition}[Zeilen- und Spaltenkomponente]
	Für $A \in \{0,1\}^{n\times n}$ und $j \in \{0,\ldots,n-1\}$ heißen
	\[
		r_j := \sum_{i=0}^{n-1} A_{ij}\cdot 2^i
		\qquad\text{und}\qquad
		c_j := j \cdot 2^n
	\]
	die \emph{Zeilenkomponente} bzw. \emph{Spaltenkomponente} der Signatur
	$\sigma_j = r_j + c_j$.
\end{definition}

\begin{definition}[Reziproke Skalierung]
	Die reziproke Skalierung der Signatur $\sigma_j$ ist definiert als
	\[
		\hat\sigma_j := \sigma_j \cdot 2^{-n} \in \mathbb{Q}_{\geq 0}.
	\]
\end{definition}

\begin{definition}[Ganzzahl- und Fraktionalanteil]
	Für $x \in \mathbb{Q}_{\geq 0}$ bezeichne $\lfloor x \rfloor$ den
	Ganzzahlanteil und $\{x\} := x - \lfloor x \rfloor \in [0,1)$ den
	Fraktionalanteil von $x$.
\end{definition}

\subsection{Grundidee}

\begin{beispiel}
	Für $n=4$ und Spalte $j=2$ mit Spaltenvektor $[1,0,1,1]^T$ gilt
	\[
		r_2 = 2^0 + 2^2 + 2^3 = 1 + 4 + 8 = 13, \qquad
		c_2 = 2 \cdot 2^4 = 32, \qquad
		\sigma_2 = 45.
	\]
	Reziproke Skalierung mit $2^{-4} = 1/16$ liefert
	\[
		\hat\sigma_2 = \frac{45}{16} = 2{,}8125
		\;=\; \underbrace{2}_{\lfloor \hat\sigma_2 \rfloor = j}
		\;+\; \underbrace{0{,}8125}_{\{\hat\sigma_2\} = r_2/2^4 = 13/16}.
	\]
	Der Ganzzahlanteil liefert exakt den Spaltenindex $j=2$ zurück, der
	Fraktionalanteil exakt die normierte Zeileninformation $r_2/2^n$ –
	und das, obwohl $2^{-4}$ als „kleine'' Zahl erscheint.
\end{beispiel}

\section{Hauptresultate}

\subsection{Exakte reziproke Zerlegung}

\begin{satz}[Reziproke Zerlegung der Signatur]
\label{satz:zerlegung}
	Für alle $n \geq 1$, alle $j \in \{0,\ldots,n-1\}$ und alle
	$A \in \{0,1\}^{n\times n}$ gilt für $\hat\sigma_j = \sigma_j \cdot 2^{-n}$:
	\[
		\lfloor \hat\sigma_j \rfloor = j
		\qquad\text{und}\qquad
		\{\hat\sigma_j\} = \frac{r_j}{2^n}.
	\]
\end{satz}

\begin{proof}
	Nach Definition gilt $\sigma_j = r_j + j\cdot 2^n$ mit
	$r_j = \sum_{i=0}^{n-1}A_{ij}2^i$. Da $A_{ij}\in\{0,1\}$ für alle $i$, gilt
	\[
		0 \;\leq\; r_j \;\leq\; \sum_{i=0}^{n-1} 2^i \;=\; 2^n - 1 \;<\; 2^n.
	\]
	Multiplikation mit $2^{-n} > 0$ liefert
	\[
		\hat\sigma_j \;=\; \sigma_j \cdot 2^{-n}
		\;=\; \frac{r_j}{2^n} + j.
	\]
	Da $0 \leq r_j/2^n < 1$ und $j \in \mathbb{Z}_{\geq 0}$, ist $j$ bereits der
	Ganzzahlanteil von $\hat\sigma_j$ und $r_j/2^n \in [0,1)$ bereits der
	Fraktionalanteil. Aus der Eindeutigkeit der Zerlegung
	$x = \lfloor x \rfloor + \{x\}$ mit $\{x\}\in[0,1)$ für $x \in \mathbb{Q}_{\geq0}$
	folgt
	\[
		\lfloor \hat\sigma_j \rfloor = j, \qquad \{\hat\sigma_j\} = \frac{r_j}{2^n}. \qedhere
	\]
\end{proof}

\begin{korollar}[Konstanz der Trennung]
\label{kor:konstanz}
	Die Breite des Wertebereichs des Fraktionalanteils, $[0,1)$, ist
	\emph{unabhängig von $n$}. Insbesondere gilt für alle $n_1 \neq n_2$ und
	beliebige zugehörige Signaturen $\sigma^{(n_1)}_j, \sigma^{(n_2)}_{j'}$:
	\[
		\{\sigma^{(n_1)}_j \cdot 2^{-n_1}\},\;
		\{\sigma^{(n_2)}_{j'} \cdot 2^{-n_2}\} \;\in\; [0,1).
	\]
\end{korollar}

\begin{proof}
	Folgt unmittelbar aus Satz~\ref{satz:zerlegung}, da $\{x\}\in[0,1)$ per
	Definition für jedes $x\in\mathbb{Q}_{\geq0}$ gilt, unabhängig von der
	Größe von $n$.
\end{proof}

\begin{bemerkung}
	Korollar~\ref{kor:konstanz} ist der formale Kern des „Paradoxons``: Obwohl
	$\sigma_j$ selbst mit $O(2^n)$ wächst (exponentiell in $n$), bleibt das
	durch $2^{-n}$ skalierte Bild $\hat\sigma_j$ in einem Intervall der
	\emph{konstanten} Länge $n$ (nämlich $[0,n)$, da $j \leq n-1$), und
	dessen Fraktionalanteil sogar in $[0,1)$ -- einer von $n$
	\emph{vollständig unabhängigen} Größe. Der scheinbare „Verlust`` durch
	das verschwindend kleine $2^{-n}$ wird also vollständig durch eine
	\emph{strukturelle Normierung} kompensiert: Information wird nicht
	skaliert, sondern \emph{aufgeteilt}.
\end{bemerkung}

\subsection{Bijektivität und Informationserhalt}

\begin{lemma}[Bijektivität der reziproken Abbildung]
\label{lem:bijektiv}
	Die Abbildung
	\[
		\Phi_n : \{0,1\}^n \times \{0,\ldots,n-1\} \longrightarrow
		\{0,\ldots,n-1\} \times \left\{\tfrac{0}{2^n},\tfrac{1}{2^n},\ldots,\tfrac{2^n-1}{2^n}\right\},
	\]
	\[
		\Phi_n(A_{\cdot j}, j) = (j,\ r_j/2^n) = (\lfloor\hat\sigma_j\rfloor, \{\hat\sigma_j\})
	\]
	ist eine Bijektion.
\end{lemma}

\begin{proof}
	\textbf{Wohldefiniertheit:} Nach Satz~\ref{satz:zerlegung} liegt
	$r_j/2^n$ stets in der angegebenen diskreten Menge von $2^n$ äquidistanten
	Werten in $[0,1)$, da $r_j \in \{0,\ldots,2^n-1\}$ ganzzahlig ist.

	\textbf{Injektivität:} Seien $(A_{\cdot j}, j) \neq (A'_{\cdot j'}, j')$.
	Fall (a) $j \neq j'$: dann unterscheiden sich die Bilder bereits in der
	ersten Komponente. Fall (b) $j = j'$, $A_{\cdot j} \neq A'_{\cdot j'}$:
	Da die Abbildung $A_{\cdot j} \mapsto r_j = \sum_i A_{ij}2^i$ die
	Standard-Binärdarstellung einer Zahl in $\{0,\ldots,2^n-1\}$ ist, und
	diese Darstellung eindeutig ist, folgt $r_j \neq r'_j$ und somit
	$r_j/2^n \neq r'_j/2^n$.

	\textbf{Surjektivität:} Für jedes Paar $(j, k/2^n)$ mit
	$k\in\{0,\ldots,2^n-1\}$ liefert die Binärdarstellung von $k$ als Vektor
	$A_{\cdot j} \in \{0,1\}^n$ ein Urbild, da jede Zahl $k \in \{0,\ldots,2^n-1\}$
	eine eindeutige Binärdarstellung mit $n$ Stellen besitzt. \qedhere
\end{proof}

\begin{satz}[Informationserhalt ohne Zusatzspeicher]
\label{satz:infoerhalt}
	Die Signatur $\sigma_j \in \{0,\ldots,n\cdot 2^n - 1\}$ kodiert das Paar
	$(j, A_{\cdot j}) \in \{0,\ldots,n-1\}\times\{0,1\}^n$ verlustfrei in
	einem einzigen Maschinenwort der Bitlänge
	\[
		B(n) = \left\lceil \log_2(n\cdot 2^n) \right\rceil
		= n + \left\lceil \log_2 n \right\rceil,
	\]
	ohne dass $j$ und $A_{\cdot j}$ in getrennten Speicherfeldern abgelegt
	werden müssen.
\end{satz}

\begin{proof}
	Nach Lemma~\ref{lem:bijektiv} ist $\Phi_n$ bijektiv, also ist die
	Gesamtabbildung
	\[
		(j, A_{\cdot j}) \;\longmapsto\; \sigma_j = r_j + j\cdot 2^n
	\]
	als Komposition von $\Phi_n$ mit der (ebenfalls bijektiven, da
	$2^n$-adische Stellenwertdarstellung) Abbildung
	$(j, r_j/2^n) \mapsto j\cdot 2^n + r_j$ injektiv. Der Bildbereich ist
	$\{0,\ldots, (n-1)\cdot 2^n + (2^n-1)\} = \{0,\ldots,n\cdot2^n-1\}$, also
	$n\cdot 2^n$ verschiedene Werte, was $\lceil\log_2(n\cdot2^n)\rceil$ Bit
	benötigt.

	Die getrennte Speicherung von $j\in\{0,\ldots,n-1\}$
	($\lceil\log_2 n\rceil$ Bit) und $A_{\cdot j}\in\{0,1\}^n$ ($n$ Bit)
	benötigt in der Summe ebenfalls $n + \lceil\log_2 n\rceil$ Bit – jedoch
	in zwei getrennten Feldern, was in realen Implementierungen wegen
	Speicherausrichtung (\emph{alignment}) und Strukturkosten
	(\emph{struct padding}) typischerweise ein zusätzliches Maschinenwort
	(z.\,B. $32$ oder $64$ Bit) erfordert. Die Signatur $\sigma_j$ vermeidet
	dies vollständig, da $j$ und $r_j$ via Satz~\ref{satz:zerlegung} aus
	einer einzigen Zahl in $O(1)$ durch die reziproke Skalierung
	$\sigma_j \cdot 2^{-n}$ (äquivalent: Ganzzahldivision durch $2^n$
	und Modulo $2^n$) zurückgewonnen werden. \qedhere
\end{proof}

\begin{korollar}[Der paradoxe Gewinn]
\label{kor:paradox}
	Der Ãœbergang von $2^n$ (Multiplikation, Kodierung) zu $2^{-n}$
	(Division, Dekodierung) ist algebraisch trivial (eine einzige
	Multiplikation mit dem Kehrwert), bewirkt aber gemäß
	Satz~\ref{satz:zerlegung} und Satz~\ref{satz:infoerhalt} die vollständige,
	verlustfreie Trennung zweier strukturell verschiedener Informationsklassen
	(Position $j$ und Inhalt $A_{\cdot j}$) bei \emph{negativen} Zusatzkosten
	gegenüber getrennter Speicherung. Der „Gewinn`` entsteht somit nicht durch
	zusätzliche Information, sondern durch \emph{Eliminierung struktureller
	Redundanz}, die bei naiver getrennter Kodierung anfallen würde.
\end{korollar}

\subsection{Konsequenz für die Laufzeit $O(n^3)$}

\begin{satz}[Reziproke Zerlegung als Voraussetzung für $O(n^3)$]
\label{satz:laufzeitkonsequenz}
	Die in Satz~\ref{satz:zerlegung} bewiesene reziproke Zerlegung ist eine
	notwendige Voraussetzung für die $O(n^2)$-Laufzeit der
	Rotation-Match-Operation (Schritt~3 des Subgraph Algorithmus, vgl.
	\cite{epp2026subgraph}), und damit für die Gesamtlaufzeit $O(n^3)$.
\end{satz}

\begin{proof}
	Im Rotation-Match-Schritt wird für jede der $n$ zyklischen Rotationen die
	\emph{Zeilenkomponente} $r_j = \sigma_j \bmod 2^n$ benötigt (vgl. Listing
	„Rotation-Match`` in \cite{epp2026subgraph}, Zeilen \texttt{row\_sigs\_A
	= [sig \% col\_weight\_A for sig in sig\_A]}). Diese Operation ist
	äquivalent zur Extraktion des Fraktionalanteils $\{\sigma_j\cdot2^{-n}\}$
	skaliert mit $2^n$, gemäß Satz~\ref{satz:zerlegung}.

	Wäre die Zerlegung \emph{nicht} in $O(1)$ pro Signatur durchführbar
	(etwa weil $j$ und $r_j$ erst durch eine Suche oder einen
	Vergleichsschritt mit Aufwand $\omega(1)$ getrennt werden müssten),
	so entstünde pro Rotation ein zusätzlicher Faktor $f(n) = \omega(1)$,
	und die Gesamtlaufzeit der Rotation-Match-Phase wäre
	\[
		O\bigl(n \cdot (n^2 + n^2)\cdot f(n)\bigr) = \omega(n^3).
	\]
	Da nach Lemma~\ref{lem:bijektiv} und Satz~\ref{satz:zerlegung} die
	Zerlegung $\sigma_j \mapsto (j, r_j)$ jedoch durch die arithmetischen
	$O(1)$-Operationen Ganzzahldivision durch $2^n$ (bzw. Multiplikation mit
	$2^{-n}$) und Modulo $2^n$ erfolgt – beides Bitschiebeoperationen
	(\texttt{shift}/\texttt{mask}), da $2^n$ eine Zweierpotenz ist – gilt
	$f(n) = O(1)$, und die in \cite{epp2026subgraph} (Satz „Laufzeit``)
	gezeigte Schranke
	\[
		T_{\text{Rotation-Match}}(n) = O(n\cdot(n^2+n^2)) = O(n^3)
	\]
	bleibt gültig. \qedhere
\end{proof}

\begin{bemerkung}
	Satz~\ref{satz:laufzeitkonsequenz} zeigt, dass der „paradoxe Gewinn``
	keine bloße ästhetische Beobachtung ist, sondern eine
	\emph{Korrektheitsvoraussetzung} der in \cite{epp2026subgraph} bewiesenen
	$\Theta(n^3)$-Optimalität: Ohne die $O(1)$-Zerlegbarkeit via $2^{\pm n}$
	würde die Gesamtschranke aus Lemma „Unabhängigkeit der Rotationsvergleiche``
	(\cite{epp2026subgraph}) um den Faktor $f(n)$ verschlechtert, was die
	scharfe untere Schranke $\Omega(n^{3-\varepsilon})$ nicht mehr matchen
	würde.
\end{bemerkung}

\section{Veranschaulichung}

\subsection{Divergenz von $2^n$ und $2^{-n}$}

Abbildung~\ref{fig:plot1} zeigt die exponentielle Divergenz von $2^n$
(Spaltengewicht, linke Achse, logarithmisch) und die exponentielle
Konvergenz von $2^{-n}$ gegen $0$ (rechte Achse, logarithmisch) für
$n=1,\ldots,16$. Beide Kurven verlaufen exakt spiegelsymmetrisch auf der
logarithmischen Skala, was die algebraische Reziprozität
$\log_2(2^n) = -\log_2(2^{-n})$ unmittelbar sichtbar macht. Auf den ersten
Blick legt dies nahe, dass $2^{-n}$ für große $n$ „bedeutungslos`` klein
und damit irrelevant wird.

\begin{figure}[htbp]
	\centering
	\includegraphics[width=0.75\textwidth]{plot1.pdf}
	\caption{Divergenz von $2^n$ und $2^{-n}$ für $n=1,\ldots,16$ (doppelt-logarithmische
	Darstellung). Die scheinbare „Bedeutungslosigkeit`` von $2^{-n}$ für große $n$
	ist der Ausgangspunkt des untersuchten Paradoxons.}
	\label{fig:plot1}
\end{figure}

\subsection{Konstante Trennung trotz wachsendem $n$}

Abbildung~\ref{fig:plot2} widerlegt diese Intuition: Für drei verschieden
große Zufallsgraphen ($n=4,8,12$) wird die reziprok skalierte Signatur
$\hat\sigma_j = \sigma_j\cdot2^{-n}$ in Ganzzahl- und Fraktionalanteil
zerlegt. Trotz exponentiell wachsender absoluter Werte von $\sigma_j$
(Größenordnung $O(n\cdot2^n)$) bleibt der Fraktionalanteil – die
Zeileninformation – für \emph{alle} $n$ exakt im Intervall $[0,1)$, und der
(reskalierte) Ganzzahlanteil identifiziert exakt den Spaltenindex $j/n$.
Dies ist die graphische Bestätigung von Satz~\ref{satz:zerlegung} und
Korollar~\ref{kor:konstanz}: Die „Auflösung`` der Trennung ist unabhängig
von $n$, obwohl der Wertebereich von $\sigma_j$ selbst mit $n$ wächst.

\begin{figure}[htbp]
	\centering
	\includegraphics[width=\textwidth]{plot2.pdf}
	\caption{Reziproke Skalierung $\sigma_j\cdot2^{-n}$ für $n=4,8,12$. Grüne Punkte:
	Fraktionalanteil (Zeileninformation $r_j/2^n$), stets in $[0,1)$. Blaue Kreuze:
	normierter Ganzzahlanteil (Spaltenindex $j/n$). Die Struktur der Trennung ist
	für alle drei Werte von $n$ identisch.}
	\label{fig:plot2}
\end{figure}

\subsection{Informationsgewinn gegenüber getrennter Kodierung}

Abbildung~\ref{fig:plot3} quantifiziert den in Satz~\ref{satz:infoerhalt} und
Korollar~\ref{kor:paradox} bewiesenen Gewinn. Die rote Kurve zeigt die
benötigten Bits pro Spalte bei \emph{getrennter} Speicherung von
Zeilenvektor und Spaltenindex inklusive eines typischen Strukturkosten-Overheads
von $32$ Bit (z.\,B. durch \emph{padding} bei Wortgrenzen-Ausrichtung in
C-Strukturen). Die grüne Kurve zeigt die Bitanzahl der kombinierten
Signatur $\sigma_j$, welche \emph{ohne} diesen Overhead auskommt, da $j$ und
$r_j$ gemäß Satz~\ref{satz:zerlegung} verlustfrei aus einer einzigen Zahl
rekonstruierbar sind. Die schraffierte Fläche entspricht den eingesparten
Bits – dem „paradoxen Gewinn`` in absoluten Speichereinheiten.

\begin{figure}[htbp]
	\centering
	\includegraphics[width=0.75\textwidth]{plot3.pdf}
	\caption{Vergleich des Speicherbedarfs pro Spalte: getrennte Kodierung
	(Zeilenvektor + Spaltenindex + Strukturkosten) vs. kombinierte Signatur
	$\sigma_j$. Die schraffierte Fläche zeigt den durch die reziproke
	Zerlegbarkeit eingesparten Speicherbedarf.}
	\label{fig:plot3}
\end{figure}

\subsection{Konsequenz für die Gesamtlaufzeit}

Abbildung~\ref{fig:plot4} stellt schließlich die in
Satz~\ref{satz:laufzeitkonsequenz} bewiesene Konsequenz dar: Links wird die
Operationsanzahl des Subgraph Algorithmus ($O(n^3)$) mit dem naiven
permutationsbasierten Ansatz ($O(n!\cdot n^2)$) für kleine $n$ verglichen
(logarithmische Skala) – der Unterschied wird bereits ab $n\approx 6$
dramatisch. Rechts ist das kubische Wachstum $n^3$ für größere $n$ separat
dargestellt. Ohne die in Abschnitt~3 bewiesene $O(1)$-Zerlegbarkeit der
Signatur via $2^{\pm n}$ wäre selbst diese kubische Schranke nicht erreichbar.

\begin{figure}[htbp]
	\centering
	\includegraphics[width=\textwidth]{plot4.pdf}
	\caption{Links: Vergleich der Operationsanzahl zwischen naivem
	Permutationsansatz $O(n!\cdot n^2)$ und Subgraph Algorithmus $O(n^3)$ für
	$n=1,\ldots,8$ (logarithmische Skala). Rechts: kubisches Wachstum $n^3$ für
	$n=1,\ldots,29$, die durch die reziproke Zerlegung garantierte
	Gesamtkomplexität.}
	\label{fig:plot4}
\end{figure}

\section{Verallgemeinerung: Das Reziproke-Trennungs-Prinzip}
\label{sec:verallgemeinerung}

Die in Abschnitt~3 für die Signaturfunktion des Subgraph Algorithmus
bewiesenen Resultate (Satz~\ref{satz:zerlegung}, Lemma~\ref{lem:bijektiv},
Satz~\ref{satz:infoerhalt}) besitzen eine gemeinsame Struktur, die über den
konkreten Anwendungsfall hinausweist. In diesem Abschnitt formulieren wir
diese Struktur als allgemeines Prinzip, leiten daraus neue
Schlussfolgerungen ab und weisen formal nach, dass ein zweites,
strukturell unabhängiges Beispiel — die \emph{Gleitkommadarstellung reeller
Zahlen} nach IEEE~754 \cite{ieee754} — demselben Prinzip genügt.

\subsection{Formulierung des allgemeinen Prinzips}

\begin{definition}[Reziprok trennbare Kodierung]
\label{def:rtk}
	Sei $X$ eine Menge von „kombinierten Objekten`` und seien $J$, $R$
	Mengen von „Skaleninformation`` bzw. „Forminformation``. Eine Familie
	von Kodierungsfunktionen
	\[
		\sigma^{(p)} : J \times R \longrightarrow \mathbb{Q}_{\geq 0},
		\qquad p \in P
	\]
	(parametrisiert durch einen \emph{Skalenparameter} $p$ aus einer
	Parametermenge $P$, z.\,B. $P=\mathbb{N}$) heißt
	\emph{reziprok trennbar}, wenn es eine Funktion $w:P\to\mathbb{R}_{>0}$
	mit $w(p)\to\infty$ für $p\to\infty$ gibt, sodass für alle $p\in P$,
	$j\in J$, $r\in R$ gilt:
	\begin{enumerate}[label=(\roman*)]
		\item $\sigma^{(p)}(j,r) = \rho(r) + \chi(j)\cdot w(p)$ für feste
		      Funktionen $\rho:R\to[0, w(p))$ und $\chi:J\to\mathbb{Z}_{\geq0}$
		      (additive Zerlegung mit Gewicht $w(p)$),
		\item die Skalierung $\hat\sigma^{(p)}(j,r) := \sigma^{(p)}(j,r)\cdot w(p)^{-1}$
		      erfüllt
		      \[
			      \lfloor \hat\sigma^{(p)}(j,r) \rfloor = \chi(j)
			      \quad\text{und}\quad
			      \{\hat\sigma^{(p)}(j,r)\} = \rho(r)/w(p),
		      \]
		\item die Abbildung $(j,r)\mapsto(\chi(j),\rho(r)/w(p))$ ist eine
		      Bijektion auf ihr Bild, und
		\item das Bild $\{\rho(r)/w(p) : r\in R\}$ ist als Teilmenge von
		      $[0,1)$ \emph{unabhängig von $p$} (bis auf Umskalierung der
		      Diskretisierung).
	\end{enumerate}
\end{definition}

\begin{satz}[Reziprokes Trennungsprinzip]
\label{satz:rtp}
	Sei $\sigma^{(p)}$ eine reziprok trennbare Kodierung im Sinne von
	Definition~\ref{def:rtk}. Dann gelten die folgenden drei Aussagen, die
	wir als \emph{paradoxen Gewinn} bezeichnen:
	\begin{enumerate}[label=(G\arabic*)]
		\item \textbf{(Verlustfreiheit trotz Skalendivergenz)} Obwohl
		      $w(p)\to\infty$ für $p\to\infty$ und damit
		      $\sigma^{(p)}(j,r)\to\infty$ für festes $j>0$, bleibt die durch
		      $w(p)^{-1}$ skalierte Größe $\hat\sigma^{(p)}(j,r)$ in einem
		      Bereich, dessen Fraktionalteil-Wertebereich $[0,1)$ \emph{nicht}
		      mit $p$ wächst (Korollar zu (iv)).
		\item \textbf{(Kostenfreie Trennung)} Die Berechnung von $\chi(j)$ und
		      $\rho(r)/w(p)$ aus $\sigma^{(p)}(j,r)$ erfordert genau eine
		      Multiplikation (mit $w(p)^{-1}$) und eine
		      Floor-/Fractional-Operation — unabhängig von $p$, also $O(1)$.
		\item \textbf{(Kompressionsgewinn)} Falls $J$ und $R$ in getrennten
		      Datenstrukturen mit einem $p$-unabhängigen Overhead $\Delta>0$
		      kodiert würden, ist die kombinierte Kodierung $\sigma^{(p)}$ für
		      hinreichend strukturierte $\rho,\chi$ um genau $\Delta$ Bit
		      bzw. Speichereinheiten günstiger.
	\end{enumerate}
\end{satz}

\begin{proof}
	\textbf{(G1):} Nach (i) gilt $\hat\sigma^{(p)}(j,r) = \rho(r)/w(p) + \chi(j)$.
	Nach Voraussetzung $\rho(r)\in[0,w(p))$, also $\rho(r)/w(p)\in[0,1)$ für
	\emph{jedes} $p$. Damit ist $\{\hat\sigma^{(p)}(j,r)\} = \rho(r)/w(p)\in[0,1)$
	für alle $p\in P$ — der Wertebereich des Fraktionalanteils ist somit
	$[0,1)$, eine von $p$ unabhängige Menge. Dies gilt, obwohl
	$\sigma^{(p)}(j,r) = \rho(r)+\chi(j)w(p) \geq \chi(j)w(p) \to \infty$ für
	$\chi(j)>0$ und $p\to\infty$. Dies ist exakt die Aussage von
	Korollar~\ref{kor:konstanz} im allgemeinen Rahmen.

	\textbf{(G2):} Nach (ii) ist
	$\chi(j) = \lfloor \sigma^{(p)}(j,r)\cdot w(p)^{-1}\rfloor$ und
	$\rho(r)/w(p) = \{\sigma^{(p)}(j,r)\cdot w(p)^{-1}\}$. Beide Größen
	werden durch dieselbe arithmetische Operationsfolge — eine
	Multiplikation mit der Konstanten $w(p)^{-1}$, gefolgt von Floor bzw.
	Subtraktion des Floor-Werts — berechnet. Diese Operationsfolge hat für
	jedes feste $p$ konstante Länge, also $O(1)$ in der Größe der Eingabe
	$\sigma^{(p)}(j,r)$.

	\textbf{(G3):} Nach (iii) ist $(j,r)\mapsto(\chi(j),\rho(r)/w(p))$
	bijektiv auf ihr Bild $X^{(p)} := \chi(J)\times(\rho(R)/w(p))$. Die
	Anzahl der zur Kodierung von $X^{(p)}$ nötigen Bits ist somit
	$\lceil\log_2|X^{(p)}|\rceil = \lceil\log_2(|\chi(J)|\cdot|\rho(R)/w(p)|)\rceil$.
	Eine getrennte Kodierung von $\chi(j)\in\chi(J)$ und
	$\rho(r)/w(p)\in\rho(R)/w(p)$ in zwei Feldern benötigt
	$\lceil\log_2|\chi(J)|\rceil + \lceil\log_2|\rho(R)/w(p)|\rceil + \Delta$
	Bit, wobei $\Delta>0$ den strukturellen Overhead (Padding,
	Wortgrenzen-Ausrichtung) bezeichnet. Da
	$\lceil\log_2(ab)\rceil \leq \lceil\log_2 a\rceil + \lceil\log_2 b\rceil$
	für alle $a,b\geq1$ (Subadditivität des Logarithmus unter Aufrundung),
	ist die kombinierte Kodierung um mindestens $\Delta$ Bit günstiger.
	\qedhere
\end{proof}

\begin{bemerkung}
\label{bem:instanz}
	Die Signaturfunktion $\sigma_j = r_j + j\cdot 2^n$ des Subgraph
	Algorithmus ist eine Instanz von Definition~\ref{def:rtk} mit
	$P=\mathbb{N}$, $p=n$, $w(n)=2^n$, $J=\{0,\ldots,n-1\}$,
	$R=\{0,1\}^n$, $\chi(j)=j$ und $\rho(A_{\cdot j}) = r_j$. Satz~\ref{satz:rtp}
	spezialisiert sich in diesem Fall exakt zu
	Satz~\ref{satz:zerlegung}, Korollar~\ref{kor:konstanz} und
	Satz~\ref{satz:infoerhalt}.
\end{bemerkung}

\subsection{Schlussfolgerung: Ein Erkennungskriterium für „paradoxe Gewinne``}

Aus Satz~\ref{satz:rtp} ergibt sich unmittelbar die folgende methodische
Schlussfolgerung, die über den Subgraph Algorithmus hinausweist:

\begin{korollar}[Erkennungskriterium]
\label{kor:erkennung}
	Immer wenn eine Kodierung zweier Informationsklassen $J$ (Skala/Position)
	und $R$ (Form/Inhalt) additiv mit einem exponentiell in einem Parameter
	$p$ wachsenden Gewicht $w(p)$ kombiniert wird,
	$\sigma^{(p)} = \rho(r) + \chi(j)\cdot w(p)$ mit $\rho(r) < w(p)$ für alle
	$p$, tritt notwendig der in Satz~\ref{satz:rtp} beschriebene
	dreifache Gewinn (G1)–(G3) auf — \emph{unabhängig davon}, wie groß
	$w(p)$ wird, und insbesondere \emph{ohne dass} die Größe von $w(p)$
	selbst jemals explizit gespeichert oder verglichen werden muss
	(es genügt, dass $w(p)^{-1}$ als Skalierungskonstante bekannt ist).
\end{korollar}

\begin{proof}
	Unmittelbare Folgerung aus Satz~\ref{satz:rtp}, da die im Korollar
	genannte Voraussetzung — additive Zerlegung mit $\rho(r)<w(p)$ — exakt
	Bedingung (i) aus Definition~\ref{def:rtk} ist, und (ii)--(iv) für
	$\chi(j)=j$, $\rho=$ identische Restklassenfunktion automatisch erfüllt
	sind, sofern $\rho$ surjektiv auf $\{0,\ldots,\lceil w(p)\rceil-1\}$
	abbildet (was für ganzzahlige $\rho(r)\in[0,w(p))$ stets der Fall ist).
\end{proof}

\begin{bemerkung}[Neue Erkenntnis]
	Korollar~\ref{kor:erkennung} liefert ein \emph{prüfbares Kriterium}: Um
	zu entscheiden, ob ein gegebenes Kodierungsschema einen „paradoxen
	Gewinn`` im Sinne dieser Arbeit aufweist, genügt es, die additive Form
	$\sigma = \rho(r)+\chi(j)\cdot w(p)$ mit der Schranke $\rho(r)<w(p)$
	nachzuweisen — alle weiteren Eigenschaften (G1)–(G3) folgen dann
	\emph{automatisch}. Dies reduziert die Suche nach weiteren Instanzen des
	Phänomens auf eine rein syntaktische Inspektion von Kodierungsformeln.
	Im folgenden Abschnitt wenden wir dieses Kriterium auf ein zweites,
	von Graphen und Adjazenzmatrizen vollständig unabhängiges Beispiel an.
\end{bemerkung}

\subsection{Zweites Beispiel: Gleitkommadarstellung nach IEEE~754}
\label{subsec:ieee754}

Wir zeigen nun, dass die \emph{normalisierte Gleitkommadarstellung} reeller
Zahlen, wie sie im Standard IEEE~754 \cite{ieee754} sowie in der klassischen
Darstellung bei Knuth \cite{knuth1997} beschrieben wird, eine reziprok
trennbare Kodierung im Sinne von Definition~\ref{def:rtk} ist — mit dem
Exponenten $E$ in der Rolle des Spaltenindex $j$ und der Mantisse $M$ in der
Rolle des Zeilenvektors $A_{\cdot j}$.

\begin{definition}[Normalisierte Gleitkommazahl]
\label{def:floatdarstellung}
	Eine positive reelle Zahl $x\in\mathbb{R}_{>0}$ besitzt für jede Basis
	$b\geq2$ eine eindeutige \emph{normalisierte Darstellung}
	\[
		x = M \cdot b^{E}, \qquad M\in[1,b),\ E\in\mathbb{Z}.
	\]
	$E$ heißt \emph{Exponent} (Skaleninformation), $M$ heißt
	\emph{Mantisse} oder \emph{Signifikand} (Forminformation).
\end{definition}

\begin{lemma}[Existenz und Eindeutigkeit der normalisierten Darstellung]
\label{lem:floatexistenz}
	Für jedes $x\in\mathbb{R}_{>0}$ und jede Basis $b\geq2$ existieren
	eindeutige $E\in\mathbb{Z}$ und $M\in[1,b)$ mit $x=M\cdot b^E$, nämlich
	\[
		E = \lfloor \log_b x \rfloor, \qquad M = x\cdot b^{-E}.
	\]
\end{lemma}

\begin{proof}
	\textbf{Existenz:} Setze $E:=\lfloor\log_b x\rfloor\in\mathbb{Z}$ (wohldefiniert,
	da $x>0$ und $\log_b$ auf $\mathbb{R}_{>0}$ definiert ist) und
	$M:=x\cdot b^{-E}$. Dann gilt
	$\log_b M = \log_b x - E = \log_b x - \lfloor\log_b x\rfloor = \{\log_b x\}\in[0,1)$,
	also $M = b^{\{\log_b x\}}\in[b^0,b^1)=[1,b)$. Damit gilt $x=M\cdot b^E$
	mit $M\in[1,b)$, $E\in\mathbb{Z}$.

	\textbf{Eindeutigkeit:} Angenommen, es gelte zugleich
	$x=M\cdot b^E=M'\cdot b^{E'}$ mit $M,M'\in[1,b)$, $E,E'\in\mathbb{Z}$
	und $(M,E)\neq(M',E')$. Ohne Einschränkung sei $E<E'$, also $E'\geq E+1$.
	Aus $M\cdot b^E=M'\cdot b^{E'}$ folgt $M = M'\cdot b^{E'-E}\geq M'\cdot b\geq b$
	(da $M'\geq1$ und $E'-E\geq1$), im Widerspruch zu $M<b$. Also $E=E'$, und
	wegen $b^E>0$ folgt unmittelbar $M=M'$. \qedhere
\end{proof}

\begin{satz}[IEEE-754-Gleitkommadarstellung als reziprok trennbare Kodierung]
\label{satz:ieee_rtk}
	Für $b=2$ ist die durch Lemma~\ref{lem:floatexistenz} gegebene
	Zerlegung $x\mapsto(E,M)$ mit $M\in[1,2)$, $E\in\mathbb{Z}$ eine
	reziprok trennbare Kodierung im Sinne von Definition~\ref{def:rtk},
	mit Skalenparameter $p:=|E|_{\max}$ (maximaler darstellbarer
	Exponentenbetrag, in IEEE~754 double durch $11$ Exponentenbits auf
	$|E|\leq1023$ begrenzt \cite{ieee754}) und Gewichtsfunktion
	$w(p)=2^{p}$.
\end{satz}

\begin{proof}
	Wir prüfen die vier Bedingungen aus Definition~\ref{def:rtk} mit
	$J:=\{-p,\ldots,p\}\ni\chi(x):=E$, $R:=[1,2)\ni\rho(x):=(M-1)\cdot 2^{p}$
	(skaliert, damit $\rho(r)\in[0,2^p)=[0,w(p))$ gilt, vgl. (i)), und
	$\sigma^{(p)}(E,M) := (M-1)\cdot2^p + (E+p)\cdot2^p$.

	\textbf{(i) Additive Zerlegung:} Nach Konstruktion ist
	$\sigma^{(p)}(E,M) = \underbrace{(M-1)\cdot2^p}_{=\rho(r)\in[0,2^p)}
	+ \underbrace{(E+p)}_{=\chi(j)\in\mathbb{Z}_{\geq0}}\cdot\underbrace{2^p}_{=w(p)}$,
	also von der Form (i) mit $\rho(r)\in[0,w(p))$, da $M-1\in[0,1)$ nach
	Lemma~\ref{lem:floatexistenz} (für $b=2$) und somit $(M-1)2^p\in[0,2^p)$.
	Die Verschiebung $E\mapsto E+p\geq0$ ist nötig, damit $\chi(j)\in\mathbb{Z}_{\geq0}$
	gilt (vgl. die analoge Rolle von $j\geq0$ in der Subgraph-Signatur);
	sie ist eine Bijektion $\{-p,\ldots,p\}\to\{0,\ldots,2p\}$ und ändert
	nichts an der nachfolgenden Argumentation.

	\textbf{(ii) Reziproke Skalierung:} Für
	$\hat\sigma^{(p)} := \sigma^{(p)}(E,M)\cdot 2^{-p}
	= (M-1) + (E+p)$ gilt analog zu Satz~\ref{satz:zerlegung}, da
	$M-1\in[0,1)$ und $E+p\in\mathbb{Z}_{\geq0}$:
	\[
		\lfloor\hat\sigma^{(p)}\rfloor = E+p, \qquad
		\{\hat\sigma^{(p)}\} = M-1.
	\]
	Damit ist $\chi(j)=\lfloor\hat\sigma^{(p)}\rfloor-p$ (Exponent) und
	$\rho(r)/w(p) = \{\hat\sigma^{(p)}\}=M-1$ (Mantissenanteil) jeweils in
	$O(1)$ aus $\hat\sigma^{(p)}$ rekonstruierbar — dies ist exakt die in
	Hardware implementierte Operation des „Auseinandernehmens`` einer
	Gleitkommazahl in Vorzeichen-, Exponenten- und Mantissenbits, die in
	jeder modernen FPU als Festverdrahtung in $O(1)$ (konstanter
	Taktzykluszahl) realisiert ist.

	\textbf{(iii) Bijektivität:} Nach Lemma~\ref{lem:floatexistenz} ist
	$x\mapsto(E,M)$ eine Bijektion $\mathbb{R}_{>0}\to\mathbb{Z}\times[1,2)$
	(eingeschränkt auf $|E|\leq p$ und auf die durch die Mantissenbreite
	$m$ diskretisierten Werte $M\in\{1,1+2^{-m},\ldots,2-2^{-m}\}$ eine
	Bijektion auf eine endliche Menge der Größe $(2p+1)\cdot2^m$). Die
	Komposition mit der affinen, bijektiven Abbildung
	$(E,M)\mapsto(E+p,(M-1))$ ist somit ebenfalls bijektiv.

	\textbf{(iv) $p$-Unabhängigkeit des Fraktionalanteil-Bildbereichs:}
	Der Bildbereich von $\rho(r)/w(p)=M-1$ ist $[0,1)$ für \emph{jedes}
	$p$ (bzw. für jede Exponentenbreite), da $M\in[1,2)$ nach
	Lemma~\ref{lem:floatexistenz} unabhängig von der konkreten Größe von
	$E$ bzw. $p$ gilt. \qedhere
\end{proof}

\begin{korollar}[Paradoxer Gewinn in der Gleitkommadarstellung]
\label{kor:floatparadox}
	Nach Satz~\ref{satz:rtp}, angewendet auf die in
	Satz~\ref{satz:ieee_rtk} nachgewiesene Instanz, gelten für die
	IEEE-754-Gleitkommadarstellung die folgenden konkreten Aussagen:
	\begin{enumerate}[label=(F\arabic*)]
		\item \textbf{(G1, Verlustfreiheit):} Obwohl der Wertebereich von
		      $x=M\cdot2^E$ für wachsende $|E|$ über $2^{1023}\approx10^{308}$
		      bzw. unter $2^{-1022}\approx10^{-308}$ hinausgeht — also um
		      \emph{Hunderte Größenordnungen} variiert —, bleibt der
		      Mantissenanteil $M-1$ stets im \emph{exakt gleichen} Intervall
		      $[0,1)$ mit \emph{exakt gleicher} Diskretisierung
		      ($2^{52}$ Stufen bei \texttt{double}), unabhängig vom Exponenten.
		\item \textbf{(G2, $O(1)$-Trennung):} Vorzeichen, Exponent und
		      Mantisse werden aus dem $64$-Bit-Wort einer
		      \texttt{double}-Zahl durch genau drei
		      Bitmasken-/Shift-Operationen extrahiert — unabhängig davon, ob
		      $|E|$ nahe $0$ oder nahe $1023$ liegt.
		\item \textbf{(G3, Kompressionsgewinn):} Eine \texttt{double}-Zahl
		      benötigt $64$ Bit ($1$ Vorzeichen-, $11$ Exponenten-, $52$
		      Mantissenbit) für die kombinierte Kodierung von $E\in\{-1022,\ldots,1023\}$
		      und $M\in[1,2)$ (diskretisiert in $2^{52}$ Stufen). Eine
		      getrennte Speicherung von $E$ (als \texttt{int16}, $16$ Bit)
		      und $M$ (als $52$-Bit-Festkommazahl, aufgerundet auf $56$ Bit
		      wegen Byte-Ausrichtung) plus Vorzeichenbit benötigte
		      $16+56+1=73$ Bit — also $\Delta=9$ Bit mehr pro Zahl als die
		      kombinierte IEEE-754-Darstellung.
	\end{enumerate}
\end{korollar}

\begin{proof}
	(F1) folgt aus (G1) (Satz~\ref{satz:rtp}) mit $w(p)=2^p$,
	$p=1023$: Der Fraktionalanteil $\{\hat\sigma^{(p)}\}=M-1\in[0,1)$ ist
	nach Definition unabhängig von $p$, während
	$\sigma^{(p)}(E,M)=(M-1)2^p+(E+p)2^p$ für $E\to1023$ Werte der
	Größenordnung $2^{1023}\cdot2^{1023}$ erreicht — eine Diskrepanz von
	mehr als $600$ Dezimalstellen zwischen der Größe von $\sigma^{(p)}$ und
	der Größe seines (immer in $[0,1)$ liegenden) Fraktionalanteils.

	(F2) folgt aus (G2): Die in Satz~\ref{satz:ieee_rtk}(ii) gezeigte
	Operationsfolge — Multiplikation mit $2^{-p}$ (äquivalent: Bitshift um
	$p$ Stellen, da $2^p$ eine Zweierpotenz ist), gefolgt von Floor und
	Subtraktion — entspricht exakt der Bitmasken-Operation
	\texttt{(bits >> 52) \& 0x7FF} (Exponent) bzw.
	\texttt{bits \& 0xFFFFFFFFFFFFF} (Mantisse) auf der $64$-Bit-Repräsentation,
	beides $O(1)$-Operationen unabhängig vom Wert von $E$.

	(F3) folgt aus (G3) durch explizite Bitzählung: $|\chi(J)|=2047$
	(Werte für $E\in\{-1022,\ldots,1024\}$, davon $2$ für Spezialwerte
	reserviert, vgl. \cite{ieee754}), kodierbar in $\lceil\log_2 2047\rceil=11$
	Bit; $|\rho(R)/w(p)|=2^{52}$, kodierbar in $52$ Bit; zusammen mit
	einem Vorzeichenbit also $1+11+52=64$ Bit für die kombinierte
	Darstellung — exakt die Größe eines \texttt{double}. Die getrennte
	Variante mit Byte-Ausrichtung ($16$-Bit-\texttt{int16} für $E$ statt
	$11$ Bit, $56$-Bit-Feld für $M$ statt $52$ Bit, plus $1$ Vorzeichenbit)
	ergibt $16+56+1=73>64$ Bit, also $\Delta=9$ Bit Mehrbedarf — in
	Ãœbereinstimmung mit der allgemeinen Aussage (G3), dass
	$\Delta>0$ für jede Kodierung mit Ausrichtungs-Overhead gilt. \qedhere
\end{proof}

\begin{bemerkung}[Tiefere Analogie und neue Erkenntnis]
\label{bem:tiefeanalogie}
	Die Analogie zwischen Subgraph-Signatur und IEEE-754-Float geht über die
	rein strukturelle Übereinstimmung hinaus: In \emph{beiden} Fällen
	kodiert die \emph{Spaltenkomponente} bzw. der \emph{Exponent} eine
	\emph{diskrete Position} in einer geordneten Menge ($j\in\{0,\ldots,n-1\}$
	bzw. $E\in\{-1022,\ldots,1023\}$), während die \emph{Zeilenkomponente}
	bzw. die \emph{Mantisse} eine \emph{kontinuierliche bzw.
	hochauflösende Forminformation} kodiert ($A_{\cdot j}\in\{0,1\}^n$ bzw.
	$M\in[1,2)$). Die neue Erkenntnis aus dieser Arbeit ist, dass dieses
	Muster — \emph{Position-mal-Skala plus beschränkte Form} — ein
	\textbf{universelles Entwurfsmuster für verlustfreie, $O(1)$-trennbare
	Kodierungen} darstellt, das unabhängig vom konkreten Anwendungsgebiet
	(Graphentheorie, numerische Datentypen) durch das in
	Korollar~\ref{kor:erkennung} formulierte syntaktische Kriterium erkannt
	werden kann. Insbesondere legt dies nahe, dass \emph{jede} Kodierung,
	die historisch unabhängig voneinander auf die Form
	$\sigma=\rho(r)+\chi(j)\cdot w(p)$ mit $\rho(r)<w(p)$ konvergiert ist
	(wie es sowohl bei der Subgraph-Signatur von 2026 als auch bei der
	IEEE-754-Norm von 1985 der Fall war), \emph{notwendigerweise} ein
	lokales Optimum bezüglich (G1)–(G3) realisiert — ein Resultat, das
	weitere, noch unentdeckte Instanzen in anderen Bereichen
	(Datenkompression, Adressierungsschemata, Zeitstempelformate) vermuten
	lässt.
\end{bemerkung}

\subsection{Veranschaulichung des zweiten Beispiels}

Abbildung~\ref{fig:plot5} veranschaulicht Satz~\ref{satz:ieee_rtk} und
Korollar~\ref{kor:floatparadox}(F1) anhand von $40$ logarithmisch breit
gestreuten Zufallswerten $x\in(0{,}001,\,10000)$, also über sieben
Größenordnungen. Der Exponent $E=\lfloor\log_2 x\rfloor$ (links) variiert
entsprechend stark und trägt die gesamte „Skaleninformation`` der Zahl $x$.
Die Mantisse $M=x\cdot2^{-E}$ (rechts) bleibt dagegen für \emph{jeden} Wert
von $x$ — unabhängig davon, ob $x\approx0{,}001$ oder $x\approx10000$ — exakt
im Intervall $[1,2)$. Diese Konstanz des Mantissenbereichs ist das exakte
Analogon zu Abbildung~\ref{fig:plot2} (konstante Trennung des
Fraktionalanteils $r_j/2^n$ in $[0,1)$ unabhängig von $n$): In beiden Fällen
wird eine über viele Größenordnungen variierende Größe durch eine reziproke
Skalierung in eine Komponente mit \emph{konstantem, von der Skala
unabhängigem Wertebereich} überführt.

\begin{figure}[htbp]
	\centering
	\includegraphics[width=\textwidth]{plot5.pdf}
	\caption{Reziproke Zerlegung $x=M\cdot2^E$ für $40$ Zufallswerte
	$x\in(0{,}001,10000)$. Links: Exponent $E$ (Skaleninformation), stark
	streuend. Rechts: Mantisse $M\in[1,2)$ (Forminformation), für alle
	Werte im selben Intervall — unabhängig von der absoluten Größe von $x$.
	Dies ist das direkte Analogon zur konstanten Trennung in
	Abbildung~\ref{fig:plot2}.}
	\label{fig:plot5}
\end{figure}

\section{Zusammenfassung}

Die vorliegende Arbeit hat den scheinbar paradoxen Sachverhalt formal
nachgewiesen, dass die reziproke Größe $2^{-n}$ in der Signaturfunktion des
Subgraph Algorithmus
\begin{itemize}
	\item trotz exponentieller Konvergenz gegen $0$ keine Information
	      vernichtet (Satz~\ref{satz:zerlegung}),
	\item eine Bijektion zwischen Spaltenvektoren und Paaren
	      $(j, r_j/2^n)$ mit von $n$ unabhängigem Wertebereich des
	      Fraktionalanteils induziert (Lemma~\ref{lem:bijektiv},
	      Korollar~\ref{kor:konstanz}),
	\item den kombinierten Speicherbedarf gegenüber getrennter Kodierung von
	      Position und Inhalt strikt reduziert (Satz~\ref{satz:infoerhalt},
	      Korollar~\ref{kor:paradox}),
	\item und eine notwendige Voraussetzung für die Erreichbarkeit der
	      $\Theta(n^3)$-Optimalität des Subgraph Algorithmus ist
	      (Satz~\ref{satz:laufzeitkonsequenz}).
\end{itemize}

Der „Gewinn im Reziproken`` besteht somit darin, dass eine algebraisch
triviale Operation – Multiplikation mit dem Kehrwert einer Zweierpotenz,
äquivalent zu einer Bitschiebeoperation – eine strukturelle Entflechtung
zweier Informationsdimensionen bewirkt, deren getrennte Kodierung in
realen Systemen stets mit Mehraufwand verbunden wäre.

\subsection{Anwendungen}

Die hier bewiesenen Eigenschaften sind relevant für:
\begin{itemize}
	\item Kompakte Hashing-Schemata, in denen Position und Inhalt eines
	      Datensatzes in einem einzigen Maschinenwort kodiert werden sollen
	\item Bitfeld-basierte Datenstrukturen in eingebetteten Systemen, in
	      denen \emph{struct padding} vermieden werden muss
	\item Die Analyse weiterer polynomialer Hash-Funktionen mit
	      Zweierpotenz-Gewichten hinsichtlich Bijektivität und
	      $O(1)$-Dekodierbarkeit
	\item Lehrmaterial zur Demonstration, dass „kleine`` reziproke Größen
	      in der Informatik strukturelle, nicht nur quantitative Bedeutung
	      haben können
\end{itemize}

\subsection{Ausblick}

Die in dieser Arbeit etablierten Resultate eröffnen eine Reihe
weiterführender Forschungsrichtungen, die im Folgenden ausführlich
dargestellt werden.

\textbf{1. Verallgemeinerung auf beliebige Basen $b > 1$.}
Die gesamte Argumentation dieser Arbeit beruht auf der Basis $2$, was sowohl
mathematisch (Binärdarstellung) als auch praktisch (Bitschiebeoperationen,
\texttt{shift}/\texttt{mask} in $O(1)$) motiviert ist. Eine naheliegende
Erweiterung besteht darin, die Signaturfunktion zu
$\sigma_j^{(b)} = \sum_{i=0}^{n-1}A_{ij}\cdot b^i + j\cdot b^n$ für beliebige
ganzzahlige Basen $b \geq 2$ zu verallgemeinern. Satz~\ref{satz:zerlegung}
und Lemma~\ref{lem:bijektiv} lassen sich strukturell unverändert auf diesen
Fall übertragen, da $0 \leq r_j \leq b^n - 1$ und somit
$\hat\sigma_j^{(b)} = \sigma_j^{(b)}\cdot b^{-n} = j + r_j/b^n$ mit
$r_j/b^n\in[0,1)$ weiterhin gilt. Die entscheidende offene Frage ist jedoch,
ob die $O(1)$-Dekodierbarkeit (Satz~\ref{satz:laufzeitkonsequenz}) für $b\neq2^k$
erhalten bleibt: Division durch $b^n$ und Modulo $b^n$ sind für allgemeine
$b$ \emph{nicht} durch Bitschiebeoperationen realisierbar, sondern erfordern
echte Ganzzahldivision, deren Hardware-Latenz typischerweise höher ist als
die von \texttt{shift}-Operationen, asymptotisch jedoch weiterhin $O(1)$
bleibt (für feste Wortbreite). Es wäre lohnend, empirisch zu untersuchen, ab
welcher Basis $b$ und welcher Knotenzahl $n$ dieser konstante, aber von Null
verschiedene Mehraufwand die Gesamtlaufzeit des Subgraph Algorithmus in der
Praxis (nicht asymptotisch) messbar beeinflusst, und ob es Basen $b$ gibt,
für die der „paradoxe Gewinn`` \emph{größer} ausfällt als für $b=2$ – etwa im
Hinblick auf eine dichtere Packung von Information bei nicht-binären
Adjazenzmatrizen (gewichtete Graphen mit Kantengewichten aus
$\{0,\ldots,b-1\}$).

\textbf{2. Mehrstufige reziproke Zerlegungen für höherdimensionale
Strukturen.} Die in dieser Arbeit behandelte Zerlegung trennt genau zwei
Informationsdimensionen (Position $j$ und Inhalt $A_{\cdot j}$) durch genau
\emph{eine} reziproke Skalierung mit $2^{-n}$. Für Tensoren höherer Ordnung
– etwa Hypergraphen mit $k$-ären Relationen oder mehrschichtige
(\emph{multilayer}) Graphen – stellt sich die Frage, ob eine
\emph{mehrstufige} Signatur der Form
\[
	\sigma_{j_1,\ldots,j_k} = r + \sum_{l=1}^{k} j_l \cdot 2^{n + \sum_{m<l} \log_2 N_m}
\]
(mit geeigneten Normierungskonstanten $N_m$) eine analoge Kette von
ineinander geschachtelten reziproken Zerlegungen
$\hat\sigma \cdot 2^{-n}, \{\hat\sigma\}\cdot N_1^{-1}, \ldots$ zulässt, bei
der jede Stufe wieder Korollar~\ref{kor:konstanz} (Konstanz der Trennung,
unabhängig von den Originalgrößen) erfüllt. Eine vollständige formale
Verallgemeinerung von Satz~\ref{satz:infoerhalt} auf $k$ Dimensionen mit
einem Beweis, dass auch hier jede Dekodierungsstufe in $O(1)$ erfolgt, steht
noch aus und wäre direkt relevant für die Erweiterung des Subgraph
Algorithmus auf gewichtete und gerichtete Graphen sowie auf Hypergraphen,
wie sie bereits als Erweiterungsmöglichkeit in \cite{epp2026subgraph}
genannt werden.

\textbf{3. Informationstheoretische Einordnung als Entropie-Zerlegung.}
Korollar~\ref{kor:paradox} wurde in dieser Arbeit operational (Bits,
Speicherplatz) formuliert. Eine tiefere informationstheoretische Fundierung
würde die Zerlegung $\sigma_j \mapsto (j, r_j)$ als Zerlegung der
gemeinsamen Entropie $H(J, R) = H(J) + H(R\mid J)$ interpretieren, wobei $J$
die Zufallsvariable des Spaltenindex und $R$ die Zufallsvariable der
Zeilenkomponente bezeichnet. Da $J$ und $R$ in der Konstruktion des
Subgraph Algorithmus stochastisch unabhängig gewählt werden können (der
Spaltenindex ist deterministisch durch die Position gegeben, der
Zeilenvektor ist die eigentliche „Nutzinformation``), gilt
$H(J,R) = H(J) + H(R)$, und die reziproke Zerlegung realisiert diese additive
Entropiezerlegung \emph{architektonisch} – nicht nur statistisch. Eine
formale Ausarbeitung dieser Sichtweise könnte den „paradoxen Gewinn`` mit dem
Begriff des \emph{kanonischen Codes} (im Sinne der Kolmogorov-Komplexität)
in Verbindung bringen: $\sigma_j$ wäre dann ein kanonischer, präfixfreier
Code für das Paar $(j, A_{\cdot j})$, und die Frage, ob dieser Code in einem
geeigneten Sinn \emph{optimal} (d.\,h. von minimaler erwarteter Länge) unter
allen $O(1)$-dekodierbaren Codes ist, wäre ein eigenständiges
Forschungsthema.

\textbf{4. Numerische Stabilität bei Gleitkomma-Implementierungen.} Diese
Arbeit hat die reziproke Skalierung $\sigma_j\cdot2^{-n}$ exakt in den
rationalen Zahlen $\mathbb{Q}$ behandelt (Satz~\ref{satz:zerlegung}). In
realen Implementierungen, insbesondere wenn Signaturen aus
Performance-Gründen als \texttt{double}-Gleitkommazahlen statt als
beliebig genaue Ganzzahlen (\texttt{BigInteger}) repräsentiert werden,
stellt sich die Frage der numerischen Stabilität: Für $n > 52$ übersteigt
$2^n$ die Mantissenbreite eines IEEE-754-\texttt{double} ($53$ Bit), sodass
$\sigma_j$ und $2^{-n}$ nicht mehr exakt darstellbar sind und die in
Satz~\ref{satz:zerlegung} bewiesene exakte Zerlegung in Rundungsfehler
übergeht. Eine systematische Fehleranalyse, die obere Schranken für den
Rundungsfehler von $\lfloor\hat\sigma_j\rfloor$ und $\{\hat\sigma_j\}$ in
Abhängigkeit von $n$ und der Mantissenbreite $p$ angibt, sowie ein
Vorschlag für eine numerisch stabile Variante (etwa durch Verwendung von
\texttt{shift}/\texttt{mask}-Operationen auf Ganzzahlen anstelle von
Gleitkomma-Multiplikation mit $2^{-n}$, wie in
Satz~\ref{satz:laufzeitkonsequenz} bereits angedeutet), wäre für
Implementierungen des Subgraph Algorithmus für sehr große Graphen
($n>10000$, vgl. \cite{epp2026subgraph}, Abschnitt „Ausblick``) von
unmittelbarer praktischer Relevanz.

\textbf{5. Empirische Validierung des Speichergewinns in realen
Systemen.} Abbildung~\ref{fig:plot3} illustriert den theoretischen
Speichergewinn unter der Annahme eines pauschalen Overheads von $32$ Bit
für getrennte Felder. Eine empirische Studie, die den tatsächlichen
Speicherverbrauch konkreter Implementierungen (z.\,B. in C/C++ mit
\texttt{struct}-Definitionen unter verschiedenen Compilern und
Architekturen, sowie in Python mit \texttt{NumPy}-Arrays vs. Tupellisten)
für $A$ und $j$ getrennt gegenüber der kombinierten Signatur $\sigma_j$
misst, würde die theoretischen Schranken aus
Satz~\ref{satz:infoerhalt} empirisch untermauern und gegebenenfalls
architekturspezifische Korrekturfaktoren liefern.

\textbf{6. Anwendung auf das ChatMe-Signatur-Chiffre-Schema.} Da der
Autor in verwandten Arbeiten ein Signatur-Chiffre-Verfahren für
verschlüsselte Kommunikation entwickelt, liegt es nahe zu untersuchen, ob
die hier bewiesene reziproke Zerlegung auch dort zur kompakten Kodierung
von Schlüssel- und Nachrichtenanteilen in einem einzigen Wert verwendet
werden kann, analog zur Trennung von Spaltenindex und Zeileninformation in
$\sigma_j$. Eine solche Übertragung würde jedoch eine sorgfältige
kryptographische Analyse erfordern, da die hier bewiesene Bijektivität
(Lemma~\ref{lem:bijektiv}) zwar Verlustfreiheit garantiert, aber keinerlei
Aussage über kryptographische Sicherheitseigenschaften wie
Kollisionsresistenz gegenüber einem Angreifer mit Kenntnis von $n$ macht –
im Gegenteil, die \emph{vollständige Determiniertheit} der Zerlegung könnte
in einem kryptographischen Kontext eine Schwäche darstellen, die in einer
gesonderten Arbeit analysiert werden müsste.

Zusammenfassend zeigt diese Arbeit, dass die scheinbar triviale
algebraische Beziehung zwischen $2^n$ und $2^{-n}$ in der Signaturfunktion
des Subgraph Algorithmus eine reichhaltige formale Struktur trägt, deren
vollständige Ausschöpfung – von allgemeinen Basen über
informationstheoretische Fundierung bis hin zu kryptographischen
Anwendungen – Gegenstand zukünftiger Forschung sein sollte.

\subsection{Implementierungshinweis}

Der Code des Subgraph Algorithmus, auf dem diese Arbeit aufbaut, ist
verfügbar unter: \url{https://github.com/hjstephan86/science}.

\begin{thebibliography}{9}

\bibitem{epp2026subgraph}
Epp, S.: \emph{Der Subgraph Algorithmus}. Eigenverlag, 2026.

\bibitem{abboud2015}
Abboud, A., Backurs, A., Vassilevska Williams, V.:
\emph{Tight Hardness Results for LCS and Other Sequence Similarity Measures}.
In: Proceedings of the 56th Annual IEEE Symposium on Foundations of
Computer Science (FOCS), 2015, S. 59--78.

\bibitem{cormen2009}
Cormen, T. H., Leiserson, C. E., Rivest, R. L., Stein, C.:
\emph{Introduction to Algorithms}. 3. Auflage, MIT Press, 2009.

\bibitem{knuth1997}
Knuth, D. E.: \emph{The Art of Computer Programming, Volume 2:
Seminumerical Algorithms}. 3. Auflage, Addison-Wesley, 1997.

\bibitem{ieee754}
IEEE Standards Association: \emph{IEEE Standard for Floating-Point
Arithmetic (IEEE 754-2019)}. 2019.

\bibitem{shannon1948}
Shannon, C. E.: \emph{A Mathematical Theory of Communication}.
Bell System Technical Journal, 27(3), 1948, S. 379--423.

\bibitem{cover2006}
Cover, T. M., Thomas, J. A.: \emph{Elements of Information Theory}.
2. Auflage, Wiley-Interscience, 2006.

\end{thebibliography}

\end{document}

